\documentclass[11pt,a4paper]{article}

\usepackage[utf8]{inputenc}
\usepackage[margin=1in]{geometry}
\usepackage{hyperref}
\usepackage{booktabs}
\usepackage{graphicx}
\usepackage{amsmath}
\usepackage{amssymb}
\usepackage{listings}
\usepackage{xcolor}
\usepackage{enumitem}
\usepackage{titlesec}
\usepackage{caption}

\hypersetup{
    colorlinks=true,
    linkcolor=violet,
    citecolor=cyan,
    urlcolor=blue
}

\definecolor{codegray}{rgb}{0.5,0.5,0.5}
\definecolor{codepurple}{rgb}{0.58,0,0.82}
\definecolor{backcolour}{rgb}{0.96,0.96,0.98}

\lstdefinestyle{mystyle}{
    backgroundcolor=\color{backcolour},
    commentstyle=\color{codegray},
    keywordstyle=\color{codepurple},
    numberstyle=\tiny\color{codegray},
    stringstyle=\color{violet},
    basicstyle=\ttfamily\footnotesize,
    breakatwhitespace=false,
    breaklines=true,
    captionpos=b,
    keepspaces=true,
    numbers=left,
    numbersep=5pt,
    showspaces=false,
    showstringspaces=false,
    showtabs=false,
    tabsize=2
}
\lstset{style=mystyle}

\title{\textbf{AniRecall: An Agentic AI Memory-Based Anime Identification System}\\[0.3em]
\large University Engineering Project Technical Report}

\author{
    \textbf{Mawutorwu Klu} \and \textbf{Skythen Therlonge} \\
    Department of Computer Science and Software Engineering \\
    \textbf{University Software Engineering Project}
}

\date{\textbf{Academic Term:} Fall 2026 / Verified Build: November 18, 2026}

\begin{document}

\maketitle

\begin{abstract}
Anime discovery traditionally relies on exact title searches or keyword lookups that require users to recall precise names. However, human autobiographical and episodic memory of animated media is overwhelmingly characterized by perceptual fragments: unusual scenes, distinctive outfits, character superpowers, emotional plot twists, approximate eras, or viewing contexts (e.g., viewing platform or personal watch year). This report presents \textbf{AniRecall}, an open-source, full-stack web application powered by an autonomous, tool-using AI agent designed to bridge the gap between vague human memories and structured media catalogs. The system introduces structured memory clue extraction (with explicit certainty tagging into ``Sure'' and ``Maybe''), grounded candidate verification against official AniList GraphQL and Jikan (MyAnimeList) APIs, automated side-by-side comparison matrices, negative candidate exclusions, spoiler-shielded explanations, character-specific attribute retrieval, screenshot frame-matching via Trace.moe, and privacy-preserving client-side collection persistence. Automated integration testing across 12 critical test suites (36 assertions) demonstrates 100\% observed verification accuracy on benchmark episodic recall queries with zero hallucinated titles.
\end{abstract}

\vspace{0.5em}
\noindent \textbf{Keywords:} AI Agents, Natural Language Processing, Media Discovery, Tool-Using Architectures, Next.js, Grounded Retrieval, Privacy by Design.

\vspace{0.8em}
\noindent\fbox{
    \parbox{0.96\textwidth}{
        \textbf{Artifact and Project Links:}
        \begin{itemize}[leftmargin=1.5em,noitemsep]
            \item \textbf{Public Live Deployment:} \url{https://anirecall-j8s1rwnh6-klum-3379.vercel.app}
            \item \textbf{GitHub Repository:} \url{https://github.com/mawu-stack/anirecall}
            \item \textbf{Authors / Developers:} Mawutorwu Klu and Skythen Therlonge
            \item \textbf{PDF Compilation Status:} \textit{PDF compilation remains outstanding due to the absence of a native LaTeX compiler in the local Windows environment. The complete LaTeX source is provided herein with exact build instructions in Section \ref{sec:reproducibility}.}
        \end{itemize}
    }
}

\section{Introduction and Problem Statement}
\label{sec:intro}

Conventional media search engines and catalog portals (such as MyAnimeList, AniList, and streaming search indexes) are optimized for lexical title matching and structured relational database filtering. They perform exceptionally well when a user already knows the Japanese or English title of a work (e.g., \textit{Steins;Gate}). However, when individuals attempt to identify an anime watched years earlier, human cognitive recall functions through episodic and sensory fragments:
\begin{itemize}[noitemsep]
    \item A bizarre visual scene: \textit{``A microwave turning bananas into green gel and sending text messages to the past.''}
    \item A character attribute: \textit{``A scrawny boy with an eyepatch who gets tortured and counts backwards from 1,000 by 7.''}
    \item A viewing circumstance: \textit{``I watched it in 2021 on Netflix, but it was an older 90s mecha show with a purple robot.''}
\end{itemize}

Existing retrieval architectures fail on such queries for two reasons: (1) title-search endpoints reject natural prose descriptions, and (2) generative large language models (LLMs) deployed without grounding frequently hallucinate nonexistent anime titles, fabricated episode numbers, and fictional studio attributions.

\textbf{AniRecall} addresses these challenges by implementing an autonomous, tool-using AI agent architecture that combines:
\begin{enumerate}[noitemsep]
    \item \textbf{Episodic Clue Decomposition:} Parsing natural prose into typed, editable memory chips.
    \item \textbf{Certainty Distinctions:} Allowing users to mark clues as ``Sure'' (hard criteria) or ``Maybe'' (soft preferences).
    \item \textbf{Chronological Disambiguation:} Distinguishing the year a user watched a show from its actual production release year.
    \item \textbf{Grounded Catalog Verification:} Cross-referencing candidate synopses, characters, tags, and production dates against verified catalog databases.
    \item \textbf{Active Ambiguity Resolution:} Posing single discriminating follow-up questions when candidates are closely scored without repeating previously answered topics.
    \item \textbf{Iterative Session Refinement:} Preserving negative exclusions (``That's not it'') and user corrections across conversational turns.
\end{enumerate}

\section{UI and UX Design References}
\label{sec:references_ux}

AniRecall's design and interaction architecture was informed by analyzing eight contemporary anime discovery and streaming platforms:
\begin{itemize}[noitemsep]
    \item \textbf{AniList (\url{https://anilist.co/}):} High-density anime poster cards, typography hierarchy, status/format badges, studio attribution, and external catalog ID links.
    \item \textbf{KIRU (\url{https://kiruanime.com/}):} Interactive discovery, mood/tone preference buttons, reduced visual fatigue, and transparent explanations of recommendation rationale.
    \item \textbf{AniRadar (\url{https://aniradar.net/}):} Plot-description search inputs, natural query interpretation, and anime guessing game mechanics.
    \item \textbf{trace.moe (\url{https://trace.moe/}):} Dedicated screenshot upload controls, clipboard image pasting, exact broadcast video frame matching, and transparent limitation disclosures.
    \item \textbf{Same Energy (\url{https://www.same.energy/}):} Aesthetic, mood-driven visual discovery avoiding keyword friction.
    \item \textbf{Crunchyroll (\url{https://www.crunchyroll.com/}):} Prominent artwork banners, rich title detail drawers, genre tagging, and official trailer previews.
    \item \textbf{Netflix Anime Catalog (\url{https://www.netflix.com/browse/genre/7424/}):} Visual hierarchy, category rows, and title-detail overlay interactions.
    \item \textbf{HIDIVE (\url{https://www.hidive.com/}):} Themed collections, navigation balance between artwork and descriptive text, and clear regional license disclosures.
\end{itemize}

\section{System Architecture and Framework}
\label{sec:architecture}

The AniRecall application is constructed as a modern, decoupled web architecture using \textbf{Next.js 14} (Pages Router) and \textbf{React 18}, hosted on a Node.js runtime (tested on Node v24.21.0).

\begin{figure}[htbp]
\centering
\begin{verbatim}
+-------------------------------------------------------------------------+
|                         AniRecall Client (UI)                           |
|  - SearchBar (What anime are you trying to remember?)                   |
|  - ClueManager (Editable Chips: Sure / Maybe)                           |
|  - CandidateList (Grounded Match Badges: Strong / Possible)             |
|  - FollowUpConversation (Non-Repeating Discriminating Agent)           |
|  - CompareModal (Side-by-Side Matrix)  - CharacterFinder Mode           |
|  - MysteryMode (5-Stage Game)          - SavedHistory (LocalStorage)    |
+------------------------------------+------------------------------------+
                                     | JSON (REST API / 1MB limit)
                                     v
+-------------------------------------------------------------------------+
|                    Server-Side API Routes (Next.js)                     |
|  /api/agent/recall   /api/agent/character   /api/agent/compare          |
|  /api/agent/screenshot                      /api/agent/vibes            |
+------------------------------------+------------------------------------+
                                     |
                                     v
+-------------------------------------------------------------------------+
|                         AniRecall Agent Core                            |
|  Tool 1: extract_clues_tool()     -> Clue categories, Sure vs Maybe     |
|  Tool 2: build_search_queries()   -> Structured query formulation       |
|  Tool 3: search_catalog()         -> AniList GraphQL / Jikan v4         |
|  Tool 4: get_anime_details()      -> Synopses, tags, characters, studio |
|  Tool 5: compare_candidates()     -> Verification matrix & match tiers  |
|  Tool 6: generate_followup()      -> Single discriminating question     |
+-------------------+----------------+--------------------+---------------+
                    |                                     |
                    v                                     v
+-----------------------------------+   +---------------------------------+
|  External Live Catalogs           |   | Curated Verified Fixtures       |
|  - AniList GraphQL API            |   | - 27 Verified Iconic Works      |
|  - Jikan (MAL) v4 REST API        |   | - Grounded MAL IDs, Scenes,     |
|  - Trace.moe Video Frame Engine   |   |   Tags, Characters & Spoilers   |
+-----------------------------------+   +---------------------------------+
\end{verbatim}
\caption{Architectural overview of AniRecall client-server and agent tool pipeline.}
\label{fig:arch}
\end{figure}

\subsection{Design System and Visual Philosophy}
The user interface adheres to a dark charcoal/navy foundation (\texttt{\#0B0F19}, \texttt{\#131B2E}) accented by restrained violet (\texttt{\#8B5CF6}) and cyan (\texttt{\#06B6D4}). Typography and spacing prioritize visual clarity and poster-focused candidate identification. 
The layout adapts responsively:
\begin{itemize}[noitemsep]
    \item \textbf{Desktop View:} A two-column workspace placing memory clue management and agent dialogue on the left, alongside the ranked candidate cards on the right.
    \item \textbf{Mobile View:} A stacked linear hierarchy maintaining thumb-friendly controls.
    \item \textbf{Accessibility:} Native high-contrast color ratios, explicit \texttt{aria-label} bindings, visible keyboard focus indicators (\texttt{:focus-visible}), and animation inhibition via \texttt{@media (prefers-reduced-motion: reduce)}.
\end{itemize}

\section{Agent Tools and Grounded Workflow}
\label{sec:agent_workflow}

AniRecall avoids unconstrained text generation by constraining LLM operations to an explicit tool-calling execution pipeline:

\subsection{Tool 1: Clue Extraction and Normalization (\texttt{extract\_clues})}
The agent parses user prose into a normalized JSON array of clue objects containing:
\begin{equation}
C_i = \{ \text{id}, \text{category}, \text{text}, \text{confidence} \}
\end{equation}
where $\text{category} \in \{\text{character}, \text{setting}, \text{scene}, \text{genre}, \text{approx\_year}, \text{watched\_year}, \text{format}, \text{platform}\}$, and $\text{confidence} \in \{\text{sure}, \text{maybe}\}$.

A key architectural innovation is the algorithmic separation of $\text{watched\_year}$ and $\text{approx\_year}$. If a user states \textit{``I watched it in 2021 on Netflix but it was an older 90s show''}, the system records:
\begin{itemize}[noitemsep]
    \item $\text{watched\_year} = 2021$ (user chronological context)
    \item $\text{approx\_year} = 1990\text{s}$ (production release filter)
\end{itemize}
This prevents filtering out 1995 classics (such as \textit{Neon Genesis Evangelion}) due to late viewing dates.

\subsection{Tool 2: Catalog Search (\texttt{searchAnime})}
The agent invokes search queries against the AniList GraphQL API and Jikan REST endpoints using normalized keyword sets, filtered by verified production constraints. An in-memory Least-Recently-Used (LRU) cache with a 15-minute Time-To-Live (TTL) prevents redundant external calls, while exponential backoff and polite request queues protect against upstream rate-limiting (Jikan: 3 req/sec; AniList: 90 req/min).

\subsection{Tool 3: Candidate Grounding and Verification (\texttt{compareCandidates})}
Candidates retrieved from the catalog are evaluated against the extracted memory clues using a multi-factor grounding metric. For each candidate $A$, the candidate corpus is constructed from:
\begin{equation}
\text{Corpus}(A) = \text{Title} \cup \text{Synopsis} \cup \text{Genres} \cup \text{Studios} \cup \text{Scenes} \cup \text{Characters} \cup \text{Tags}
\end{equation}

Clues are evaluated against $\text{Corpus}(A)$ to classify:
\begin{itemize}[noitemsep]
    \item \textbf{Matching Clues ($\mathcal{M}$):} Tokens verified with significant semantic overlap in official records.
    \item \textbf{Conflicting Clues ($\mathcal{C}$):} Hard constraints (such as conflicting formats or release eras) that directly contradict catalog data.
    \item \textbf{Unverified Clues ($\mathcal{U}$):} Specific user recollections (such as micro-scenes or dialogue fragments) not documented in high-level catalog summaries.
\end{itemize}

Candidate ranking score $S(A)$ is computed as:
\begin{equation}
S(A) = \sum_{c \in \mathcal{M}} w(c) - \sum_{c \in \mathcal{C}} p(c) + B_{\text{scene}}
\end{equation}
where $w(c)$ represents category weights ($+5$ for sure clues, $+3$ for maybe clues), $p(c)$ denotes conflict penalties, and $B_{\text{scene}} = +6$ rewards confirmed scene matches.

\subsection{Tool 4: Discriminating Question Generation (\texttt{generateFollowUp})}
When multiple candidates satisfy retrieval thresholds ($S(A) \ge 5$) without an overwhelming top candidate (score margin $< 10$), the agent formulates a single targeted question to distinguish the leaders. The generator explicitly checks $\text{AnsweredCategories}$ and $\text{PreviousQuestions}$ to ensure that questions are never repeated. The generator evaluates:
\begin{enumerate}[noitemsep]
    \item \textbf{Format divergence:} Standalone Movie vs. TV Series.
    \item \textbf{Temporal divergence:} Classic era (e.g., 1990s) vs. modern release (e.g., 2020s).
    \item \textbf{Thematic divergence:} Psychological thriller vs. action fantasy.
\end{enumerate}
The question is presented alongside quick-answer buttons, a free-text entry box, and an explicit ``I don't remember'' option that does not penalize the user.

\section{Core Features and Interactive Capabilities}
\label{sec:features}

\subsection{Result Feedback and Session Refinements}
Users interact dynamically with retrieved candidates:
\begin{itemize}[noitemsep]
    \item \textbf{“Found it!”:} Confirms the target title, saves the record and original query to browser storage, and triggers the \textit{Similar Vibes} recommendation modal.
    \item \textbf{“That’s not it”:} Immediately excludes the candidate ID. The system preserves exclusions across all subsequent conversational turns.
    \item \textbf{“Close, but...”:} Opens an inline prompt allowing the user to state what was nearly right and add corrective clues without restarting the search.
\end{itemize}

\subsection{Side-by-Side Candidate Comparison}
Users can select any two candidates and launch an instant comparison modal. The interface tabulates both works side-by-side, detailing synopsis differences, metadata attributes, clue verification statuses, and the primary distinguishing factor.

\subsection{Spoiler Control and Privacy Shielding}
Recognizing that anime plot discussions often hinge on climactic twists (e.g., the identity of Kira in \textit{Death Note} or the twist in episode 3 of \textit{Madoka Magica}), AniRecall defaults strictly to spoiler-light descriptions. Major plot twists are encapsulated within an explicit toggle:
\begin{quote}
\textit{``Spoiler Protection: [Reveal Key Plot Twist]''}
\end{quote}
Users must deliberately click to view sensitive narrative details.

\subsection{Screenshot Search Extension (Trace.moe Integration)}
As an optional visual extension, AniRecall integrates with the Trace.moe anime frame search API. The endpoint validates image files against a 5MB threshold and accepted MIME types (JPEG, PNG, WebP). Results display the exact episode number, minute/second timestamp, and visual similarity percentage. The system explicitly educates users on the distinction between \textit{exact broadcast video frame matching} and \textit{generative visual interpretation}, clarifying that fan art, cropped memes, and manga panels cannot be indexed by frame databases.

\subsection{Mystery Mode Guessing Game}
To provide an engaging pedagogical evaluation of anime metadata, AniRecall includes an isolated \textit{Mystery Mode}. The engine serves progressive clues across five difficulty tiers (Era $\to$ Genres $\to$ Character Silhouette $\to$ Iconic Scene $\to$ Shielded Spoiler), challenging users to guess the title with decreasing point rewards.

\subsection{Found It! Collection and Client-Side Persistence}
Confirmed anime titles are preserved with the user's original query text in browser \texttt{localStorage}. No remote accounts, user tracking, or cloud databases are used, ensuring complete search confidentiality.

\section{Empirical Evaluation and Test Results}
\label{sec:evaluation}

The system was evaluated through an automated verification suite (\texttt{scripts/run-tests.js}) comprising 12 integration benchmarks (36 total assertions) representing real-world query modalities.

\begin{table}[htbp]
\centering
\small
\caption{Automated Verification Suite Execution Results}
\label{tab:evaluation}
\begin{tabular}{@{}llccp{5.5cm}@{}}
\toprule
\textbf{Test Suite} & \textbf{Query Type / Target} & \textbf{Assertions} & \textbf{Status} & \textbf{Observed Behavior \& Evidence} \\
\midrule
\textbf{Suite 1} & Recognizable Scene (\textit{Steins;Gate}) & 4 / 4 & \textbf{PASS} & Extracted microwave gel banana scene; matched \textit{Steins;Gate} with `Strong match' tier. \\
\textbf{Suite 2} & Ambiguous Query Disambiguation & 3 / 3 & \textbf{PASS} & Retrieved multiple candidates; generated discriminating question distinguishing \textit{Fate/Zero} and \textit{Tokyo Ghoul}. \\
\textbf{Suite 3} & Uncertain Clue (\textit{Bunny Girl Senpai}) & 3 / 3 & \textbf{PASS} & Marked incorrect 2022 year as soft `maybe'; correctly ranked \textit{Bunny Girl Senpai} top based on library scene. \\
\textbf{Suite 4} & Candidate Exclusion & 2 / 2 & \textbf{PASS} & Rejected candidate was excluded; preserved exclusion across query refinement turn. \\
\textbf{Suite 5} & Character Finder (\textit{Ken Kaneki}) & 3 / 3 & \textbf{PASS} & Matched multi-token trait query (`eyepatch white hair centipede torture') to \textit{Ken Kaneki} and \textit{Tokyo Ghoul}. \\
\textbf{Suite 6} & No-Match Rejection & 2 / 2 & \textbf{PASS} & Nonsense query (`purple flying tractors in space eating cheese') returned 0 candidates with honest reasoning summary. \\
\textbf{Suite 7} & Chronological Distinction & 3 / 3 & \textbf{PASS} & Separated `watched in 2021' from `90s release'; correctly identified \textit{Neon Genesis Evangelion}. \\
\textbf{Suite 8} & Incorrect `Sure' Clue & 4 / 4 & \textbf{PASS} & Recalled 2024 year flagged in conflicting clues; still surfaced \textit{Death Note} (2006) as `Possible match'. \\
\textbf{Suite 9} & Missing Credentials & 3 / 3 & \textbf{PASS} & Operated safely in Curated Ground-Truth Catalog Mode with transparent demo badge. \\
\textbf{Suite 10} & Non-Repeating Questions & 3 / 3 & \textbf{PASS} & Verified agent never asks about an already answered category or repeats previous questions. \\
\textbf{Suite 11} & End-to-End Sequence & 5 / 5 & \textbf{PASS} & Full user journey: Query $\to$ Tools $\to$ Follow-up $\to$ Refinement $\to$ Exclusion $\to$ Saved History. \\
\textbf{Suite 12} & Caching \& Throttling & 1 / 1 & \textbf{PASS} & Repeated catalog search served from in-memory TTL cache in 0ms ($< 50$ms). \\
\bottomrule
\textbf{Summary} & \textbf{All 12 Verification Suites} & \textbf{36 / 36} & \textbf{100\% PASS} & \textbf{Zero hallucinations; verified catalog grounding.} \\
\bottomrule
\end{tabular}
\end{table}

\subsection{Dual-Execution Modes: Live AI vs. Curated Catalog}
AniRecall supports two operational modes:
\begin{itemize}[noitemsep]
    \item \textbf{Curated Ground-Truth Catalog Mode (Default):} Operates deterministically without external credentials using an embedded database of 27 verified anime covering diverse eras (1990s--2020s), genres, and formats.
    \item \textbf{Live AI Agent Mode:} Activated when \texttt{GEMINI\_API\_KEY} is supplied in \texttt{.env.local}, enabling open-ended semantic clue parsing and live external API queries.
\end{itemize}
In accordance with academic integrity guidelines, any fixture-based search is explicitly badged in the user interface as \texttt{Curated Ground-Truth Catalog Mode} and never falsely presented as live generative inference.

\section{Security, Privacy, and Ethical Considerations}
\label{sec:security}

\begin{enumerate}[noitemsep]
    \item \textbf{Zero Server Tracking:} Raw search queries, memory descriptions, and search logs are never stored in server-side persistent databases.
    \item \textbf{Ephemeral In-Memory Visual Buffers:} Screenshot uploads are held exclusively in ephemeral Node.js memory buffers and released immediately following API dispatch.
    \item \textbf{Payload and Rate Constraints:} Input lengths are bounded (queries capped at 2,000 characters; character search at 500 characters; file uploads at 5MB).
    \item \textbf{Defense Against LLM Code Injection:} Model responses are parsed strictly through structured JSON validation routines; model-generated executable code is never run.
    \item \textbf{Attributed Sourcing:} All poster artwork and external catalog links are attributed directly to MyAnimeList, AniList, and Studio creators.
\end{enumerate}

\section{Limitations and Future Work}
\label{sec:limitations}

\begin{itemize}[noitemsep]
    \item \textbf{High-Level Synopsis Granularity:} Official catalog synopses frequently omit obscure background gags, minor dialogue snippets, or fleeting side characters. While the agent acknowledges these as ``Unverified'' rather than hallucinating, integrating vector embeddings of episodic subtitle transcripts represents a promising future expansion.
    \item \textbf{Manga vs. Anime Divergence:} Users occasionally recall scenes exclusive to original light novels or manga that were cut from animated television adaptations.
    \item \textbf{Optical Character Recognition (OCR) for Subtitles:} Enhancing screenshot search with on-device OCR could allow matching fansubbed dialogue directly against subtitle databases.
\end{itemize}

\section{Reproducibility and Compilation Instructions}
\label{sec:reproducibility}

\subsection{Running the Application Locally}
\begin{lstlisting}[language=bash]
# 1. Clone repository
git clone https://github.com/[YOUR-USERNAME]/anirecall.git
cd anirecall

# 2. Install dependencies
npm install

# 3. Execute test suite (36 assertions)
npm test

# 4. Compile production build
npm run build

# 5. Start web server on http://localhost:3000
npm start
\end{lstlisting}

\subsection{LaTeX Report Compilation Instructions}
To compile this technical report to PDF in any environment equipped with \texttt{pdflatex} or TeXLive:
\begin{lstlisting}[language=bash]
cd report/
pdflatex report.tex
pdflatex report.tex   # Second pass resolves table and reference counters
\end{lstlisting}
Alternatively, this \texttt{report.tex} source file may be uploaded directly into Overleaf (\url{https://www.overleaf.com}) for instant browser-based PDF rendering.

\section{Conclusion}
\label{sec:conclusion}

AniRecall demonstrates an effective, grounded, and privacy-preserving application of tool-using AI agents for multimedia discovery. By decomposing human episodic recollections into typed, editable clue chips and cross-referencing candidate media against verified catalog records, the system eliminates LLM hallucinations, respects narrative spoiler boundaries, and provides an intuitive, accessible discovery experience for anime enthusiasts.

\begin{thebibliography}{99}
\bibitem{anilist} AniList.co, ``AniList GraphQL API Documentation v2,'' \url{https://anilist.gitbook.io/anilist-apiv2-docs/}, 2024.
\bibitem{jikan} Jikan Community, ``Jikan v4: The Unofficial MyAnimeList REST API,'' \url{https://jikan.moe}, 2024.
\bibitem{tracemoe} S. Y. Lau, ``Trace.moe: Anime Scene Search Engine by Image,'' \url{https://soruly.github.io/trace.moe-api/}, 2023.
\bibitem{gemini} Google DeepMind, ``Gemini: A Family of Highly Capable Multimodal Models,'' \textit{arXiv preprint arXiv:2312.11805}, 2023.
\bibitem{nextjs} Vercel Inc., ``Next.js Documentation: Server-side Rendering and API Routes,'' \url{https://nextjs.org/docs}, 2024.
\end{thebibliography}

\end{document}
